% from ega3-6-fr.tex
\oldpage[III]{156}
\begin{env}[6.6.5]
\emph{Suites spectrales $(c)$ et $(d)$.} Par définition, le complexe
$\mathrm H_n^{\mathrm{II}}(L_{\bullet\bullet}^{(i)})$ a pour terme de degré
$h$ le $A$-module
$C^{-h}(\mathfrak U^{(i)},\mathscr H_n(\mathscr P_\bullet^{(i)}))$, en vertu
de l'exactitude du foncteur $C^{-h}$. On a donc, par définition de l'hypertor
de deux Modules relatif à deux recouvrements (6.6.2)
\[
  {}^{(c)}\mathscr E_{pq}^2
  =\bigoplus_{q_1+q_2=q}
   \mathbf{Tor}_p^S
   \bigl(\mathfrak U^{(1)},\mathfrak U^{(2)};
         \mathscr H_{q_1}(\mathscr P_\bullet^{(1)}),
         \mathscr H_{q_2}(\mathscr P_\bullet^{(2)})\bigr).
\]
Enfin, par définition,
$\mathbf{Tor}_q^{A,\mathrm I}
(L_{\bullet\bullet}^{(1)},L_{\bullet\bullet}^{(2)})$ est un bicomplexe dont
le terme de degré $(h,k)$ est le $A$-module
$\mathbf{Tor}_q^S
(\mathfrak U^{(1)},\mathfrak U^{(2)};
 \mathscr P_h^{(1)},\mathscr P_k^{(2)})$. On a donc
\[
  {}^{(d)}\mathscr E_{pq}^2
  =\mathrm H_p\bigl(
    \mathbf{Tor}_q^S
    (\mathfrak U^{(1)},\mathfrak U^{(2)};
     \mathscr P_\bullet^{(1)},\mathscr P_\bullet^{(2)})\bigr).
\]
\end{env}

\begin{env}[6.6.6]
La théorie de l'hyperhomologie des foncteurs de bicomplexes
($0_{\mathrm I}$, 11.7.3) montre, comme dans (6.3.4), que, pour toute
résolution \emph{plate} de Cartan-Eilenberg $M_{\bullet\bullet}^{(i)}$ de
$L_{\bullet\bullet}^{(i)}$ (dans la catégorie des complexes de modules
limités inférieurement) ($i=1,2$), on a des isomorphismes canoniques de
bi-$\partial$-foncteurs
\[
\begin{aligned}
  \mathbf{Tor}_\bullet^S
  (\mathfrak U^{(1)},\mathfrak U^{(2)};
   \mathscr P_\bullet^{(1)},\mathscr P_\bullet^{(2)})
  &\xrightarrow{\sim}
   \mathrm H_\bullet
   (M_{\bullet\bullet}^{(1)}\otimes_A M_{\bullet\bullet}^{(2)})\\
  &\xrightarrow{\sim}
   \mathrm H_\bullet
   (M_{\bullet\bullet}^{(1)}\otimes_A L_{\bullet\bullet}^{(2)})
   \xrightarrow{\sim}
   \mathrm H_\bullet
   (L_{\bullet\bullet}^{(1)}\otimes_A M_{\bullet\bullet}^{(2)}).
\end{aligned}
\tag{6.6.6.1}\label{III.6.6.6.1.fr}
\]
\end{env}

\begin{env}[6.6.7]
Nous allons maintenant montrer que l'hypertor global défini dans (6.6.2), et
les six suites spectrales correspondantes, ne dépendent pas des recouvrements
ouverts affines finis $\mathfrak U^{(i)}$ qui ont servi à les définir (à des
isomorphismes canoniques près). Il suffira pour cela de montrer que si
$\mathfrak V^{(i)}$ sont deux autres recouvrements de même nature, tels que
$\mathfrak V^{(i)}$ soit \emph{plus fin} que $\mathfrak U^{(i)}$ pour
$i=1,2$, alors on a des \emph{isomorphismes canoniques} de foncteurs spectraux
\[
  {}^{(t)}\mathscr E
  (\mathfrak U^{(1)},\mathfrak U^{(2)};
   \mathscr P_\bullet^{(1)},\mathscr P_\bullet^{(2)})
  \xrightarrow{\sim}
  {}^{(t)}\mathscr E
  (\mathfrak V^{(1)},\mathfrak V^{(2)};
   \mathscr P_\bullet^{(1)},\mathscr P_\bullet^{(2)})
  \tag{6.6.7, $t$}\label{III.6.6.7.t.fr}
\]
où $t$ est remplacé par $a$, $b$, $a'$, $b'$, $c$ ou $d$.

Or, on a pour $i=1,2$ des homomorphismes de bicomplexes
\[
  C^\bullet(\mathfrak U^{(i)},\mathscr P_\bullet^{(i)})
  \longrightarrow
  C^\bullet(\mathfrak V^{(i)},\mathscr P_\bullet^{(i)})
\]
bien définis à homotopies près (G, II, 5.7.1); il en résulte déjà des
homomorphismes (6.6.7, $t$) canoniquement définis et compatibles avec les
opérateurs bords dans les aboutissements ($0_{\mathrm I}$, 11.3.2). En outre,
le calcul des termes $E_2$ des suites spectrales $(a)$, $(b)$, $(a')$, $(b')$,
montre que pour ces suites spectrales l'homomorphisme (6.6.7, $t$) est un
isomorphisme sur les termes $E_2$; comme ces suites spectrales sont
birégulières, on voit que (6.6.7, $t$) est un isomorphisme pour ces trois
foncteurs spectraux, donc un \emph{isomorphisme de bi-$\partial$-foncteurs}
pour leur aboutissement commun ($0_{\mathrm I}$, 11.1.5).

En particulier, pour des $\mathscr O_{X^{(i)}}$-Modules quasi-cohérents
$\mathscr F^{(i)}$ ($i=1,2$), l'homomorphisme canonique
\[
  \mathbf{Tor}_\bullet^S
  (\mathfrak U^{(1)},\mathfrak U^{(2)};
   \mathscr F^{(1)},\mathscr F^{(2)})
  \longrightarrow
  \mathbf{Tor}_\bullet^S
  (\mathfrak V^{(1)},\mathfrak V^{(2)};
   \mathscr F^{(1)},\mathscr F^{(2)})
\]
est bijectif; vu le calcul de (6.6.5), on voit que (6.6.7, $t$) est aussi un
isomorphisme des termes $E_2$ pour $t=c$ et $t=d$. On conclut comme ci-dessus
que (6.6.7, $t$) est aussi un isomorphisme de suites spectrales pour $t=c$ et
$t=d$.
